\documentclass[12pt,a4paper]{article}
\usepackage[utf8]{inputenc}
\usepackage[T1]{fontenc}
\usepackage[ngerman]{babel}
\usepackage{geometry}
\geometry{left=2.5cm, right=2.5cm, top=2.5cm, bottom=2.5cm}
\usepackage{amsmath, amssymb, amsthm}
\usepackage{tikz}
\usetikzlibrary{arrows.meta, decorations.pathmorphing, patterns, calc, 3d, perspective}
\usepackage{pgfplots}
\pgfplotsset{compat=1.18}
\usepackage{booktabs}
\usepackage{array}
\usepackage{xcolor}
\usepackage{hyperref}
\usepackage{float}
\usepackage{caption}
\usepackage{subcaption}
\usepackage{enumitem}
\usepackage{fancyhdr}
\usepackage{titlesec}

\definecolor{holzbraun}{RGB}{139, 90, 43}
\definecolor{hellbraun}{RGB}{210, 170, 110}
\definecolor{dunkelgrau}{RGB}{60,60,60}
\definecolor{hellgrau}{RGB}{220,220,220}
\definecolor{akzentblau}{RGB}{30,80,160}
\definecolor{warnung}{RGB}{200,60,60}
\definecolor{ok}{RGB}{40,140,40}

\pagestyle{fancy}
\fancyhf{}
\fancyhead[L]{\textbf{Holzbett 2\,m $\times$ 2\,m – Konstruktionsdokumentation}}
\fancyhead[R]{\thepage}
\fancyfoot[C]{\small Statiknachweis nach DIN EN 1995-1-1 (Holzbau)}

\theoremstyle{definition}
\newtheorem{definition}{Definition}
\theoremstyle{plain}
\newtheorem{theorem}{Satz}
\newtheorem{lemma}[theorem]{Lemma}
\theoremstyle{remark}
\newtheorem{bemerkung}{Bemerkung}

\title{
  \vspace{-1cm}
  {\LARGE \textbf{Holzbett für zwei 1\,m $\times$ 2\,m Matratzen}}\\[0.5em]
  {\large Konstruktionsplan, Aufbauanleitung und Statiknachweis}\\[0.3em]
  {\normalsize Mittelträger ohne Bodenstütze – Zwei nebeneinander liegende Lattenroste}
}
\author{Eigenbauprojekt -- Holzkonstruktion}
\date{\today}

\begin{document}

\maketitle
\thispagestyle{fancy}

\tableofcontents
\newpage

%===========================================================
\section{Übersicht und Konzept}
%===========================================================

Das Bettgestell ist für zwei nebeneinanderliegende Matratzen der Maße
$1{,}00\,\text{m} \times 2{,}00\,\text{m}$ ausgelegt. Die Gesamtliegefläche
beträgt $2{,}00\,\text{m} \times 2{,}00\,\text{m}$.

\medskip
\textbf{Kernidee:} Anstelle einer Mittelstütze zum Boden wird ein horizontaler
\emph{Längsträger} (Zentralträger) eingesetzt, der ausschließlich an Kopf- und
Fußholm aufgelagert ist. Beide Lattenroste liegen außen auf den Seitenholmen
und innen auf dem Zentralträger auf. Es gibt \textbf{keine einzige
Bodenstütze} im Bereich der Liegefläche.

\subsection{Koordinatensystem und Grundmaße}

\begin{center}
\begin{tikzpicture}[scale=1.0]
  % Grundriss Matratzen
  \fill[hellbraun!40] (0,0) rectangle (4,4);
  \fill[hellbraun!40] (4.1,0) rectangle (8.1,4);
  \draw[thick, holzbraun] (0,0) rectangle (4,4);
  \draw[thick, holzbraun] (4.1,0) rectangle (8.1,4);
  \node at (2,2) {\textbf{Matratze L}};
  \node at (2,1.5) {$1{,}0\,\text{m} \times 2{,}0\,\text{m}$};
  \node at (6.05,2) {\textbf{Matratze R}};
  \node at (6.05,1.5) {$1{,}0\,\text{m} \times 2{,}0\,\text{m}$};

  % Rahmen
  \draw[very thick, holzbraun!70!black, fill=holzbraun!20]
    (-0.4,-0.4) rectangle (8.5,4.4);

  % Mittelträger
  \fill[holzbraun!80] (3.9,-0.4) rectangle (4.2,4.4);
  \node[rotate=90, font=\small\bfseries] at (4.05,2) {Zentralträger};

  % Maßlinien
  \draw[<->, akzentblau, thick] (-0.4,-0.9) -- (8.5,-0.9)
    node[midway, below] {$\approx 205\,\text{cm}$ (Außenmaß)};
  \draw[<->, akzentblau, thick] (9.0,0) -- (9.0,4)
    node[midway, right] {$200\,\text{cm}$};
  \draw[<->, akzentblau, thick] (-1.0,0) -- (-1.0,4)
    node[midway, left] {$200\,\text{cm}$};

  % Eckbeine
  \fill[dunkelgrau] (-0.4,-0.4) rectangle (0,0);
  \fill[dunkelgrau] (8.1,-0.4) rectangle (8.5,0);
  \fill[dunkelgrau] (-0.4,4.0) rectangle (0,4.4);
  \fill[dunkelgrau] (8.1,4.0) rectangle (8.5,4.4);

  \node[font=\footnotesize] at (-0.2,-0.7) {Bein};
  \node[font=\footnotesize] at (8.3,-0.7) {Bein};

  % Legende
  \fill[holzbraun!80] (0,-1.8) rectangle (0.4,-1.5);
  \node[right, font=\small] at (0.5,-1.65) {Zentralträger};
  \fill[dunkelgrau] (3,-1.8) rectangle (3.4,-1.5);
  \node[right, font=\small] at (3.5,-1.65) {Eckbeine};
  \fill[holzbraun!20] (6,-1.8) rectangle (6.4,-1.5);
  \draw (6,-1.8) rectangle (6.4,-1.5);
  \node[right, font=\small] at (6.5,-1.65) {Rahmenholme};
\end{tikzpicture}
\captionof{figure}{Grundriss des Bettgestells (Draufsicht, schematisch)}
\end{center}

%===========================================================
\section{Stückliste und Materialien}
%===========================================================

\subsection{Bauteilliste}

\begin{table}[H]
\centering
\begin{tabular}{llccc}
\toprule
\textbf{Bauteil} & \textbf{Querschnitt} & \textbf{Anzahl} & \textbf{Länge} & \textbf{Holzart}\\
\midrule
Eckbeine           & $80 \times 80\,\text{mm}$  & 4 & $350\,\text{mm}$ & KVH Fichte \\
Seitenholme (lang) & $80 \times 60\,\text{mm}$  & 2 & $2050\,\text{mm}$ & KVH Fichte \\
Kopf-/Fußholm      & $80 \times 60\,\text{mm}$  & 2 & $2050\,\text{mm}$ & KVH Fichte \\
\textbf{Zentralträger}  & $\mathbf{100 \times 80\,\text{mm}}$ & \textbf{1} & $\mathbf{2000\,\text{mm}}$ & \textbf{KVH Fichte} \\
Auflagestreifen    & $40 \times 30\,\text{mm}$  & 4 & $2000\,\text{mm}$ & KVH Fichte \\
\bottomrule
\end{tabular}
\caption{Stückliste aller Holzbauteile (KVH = Konstruktionsvollholz, getrocknet)}
\end{table}

\subsection{Verbindungsmittel und Hilfsmaterialien}

\begin{table}[H]
\centering
\begin{tabular}{lcc}
\toprule
\textbf{Material} & \textbf{Menge} & \textbf{Verwendung}\\
\midrule
Holzschrauben $4 \times 60\,\text{mm}$ (Torx)  & 80 Stk. & Auflagestreifen, Verbindungen \\
Holzschrauben $6 \times 120\,\text{mm}$ (Torx) & 16 Stk. & Beine $\to$ Holmes \\
M8-Gewindestangen $\times$ 120\,mm + Muttern     &  8 Stk. & Kopf-/Fußholm lösbar \\
Holzdübel $\varnothing 10\,\text{mm}$            & 16 Stk. & Ausrichtung Eckverbindungen\\
Holzleim (D3, wasserfest)                        & 500\,g   & Leimverbindungen \\
Schleifpapier K80/K120/K180                      & je 2 Bögen & Oberflächenbehandlung \\
Holzöl oder Hartwachsöl                         & 1\,l     & Oberflächenschutz \\
\bottomrule
\end{tabular}
\caption{Verbindungsmittel und Hilfsmaterialien}
\end{table}

%===========================================================
\section{Technische Zeichnungen (TikZ)}
%===========================================================

\subsection{Seitenansicht – Querschnitt durch den Zentralträger}

\begin{figure}[H]
\centering
\begin{tikzpicture}[scale=1.4]
  % Boden (gestrichelt)
  \draw[dashed, gray] (-1,0) -- (9,0) node[right, gray] {Boden};

  % Linkes Bein
  \fill[dunkelgrau!80] (0,0) rectangle (0.6,2.5);
  \draw[thick] (0,0) rectangle (0.6,2.5);
  \node[font=\tiny, rotate=90] at (0.3,1.25) {Bein 80×80};

  % Rechtes Bein
  \fill[dunkelgrau!80] (7.4,0) rectangle (8.0,2.5);
  \draw[thick] (7.4,0) rectangle (8.0,2.5);

  % Linker Seitenholm
  \fill[holzbraun!60] (0.6,2.5) rectangle (7.4,3.1);
  \draw[thick] (0.6,2.5) rectangle (7.4,3.1);
  \node[font=\small] at (4.0,2.8) {Seitenholm $80\times60$\,mm};

  % Linker Auflagestreifen
  \fill[hellbraun!90] (0.6,3.1) rectangle (7.4,3.4);
  \draw[thick] (0.6,3.1) rectangle (7.4,3.4);
  \node[font=\tiny] at (2.0,3.25) {Auflage $40\times30$};

  % Lattenrost links (schematisch)
  \foreach \x in {1.0, 1.5, 2.0, 2.5, 3.0, 3.5} {
    \fill[hellgrau] (\x, 3.4) rectangle (\x+0.3, 3.6);
    \draw (\x, 3.4) rectangle (\x+0.3, 3.6);
  }
  \node[font=\tiny] at (2.5,3.75) {Lattenrost links};

  % Zentralträger
  \fill[holzbraun!90] (3.6,2.2) rectangle (4.4,3.4);
  \draw[very thick, holzbraun!30!black] (3.6,2.2) rectangle (4.4,3.4);
  \node[font=\tiny\bfseries, rotate=90] at (4.0,2.8) {Zentral};

  % Einlassung im Kopfholm
  \draw[<->, warnung, thick] (3.6,2.2) -- (3.6,2.5)
    node[left, font=\tiny, warnung] {Zapfen $30\,\text{mm}$};

  % Bemaßung Höhe
  \draw[<->, akzentblau] (-0.5,0) -- (-0.5,2.5)
    node[midway, left, font=\small] {$350\,\text{mm}$};
  \draw[<->, akzentblau] (-0.5,2.5) -- (-0.5,3.1)
    node[midway, left, font=\tiny] {$60\,\text{mm}$};
  \draw[<->, akzentblau] (-0.5,3.1) -- (-0.5,3.4)
    node[midway, left, font=\tiny] {$30\,\text{mm}$};

  % Gesamthöhe
  \draw[<->, akzentblau!50] (8.5,0) -- (8.5,3.4)
    node[midway, right, font=\small] {$\approx 440\,\text{mm}$};

  % Zapfenlager-Detail
  \draw[warnung, very thick, ->] (4.4,2.35) -- (5.5,1.8)
    node[right, font=\tiny, warnung, text width=2.5cm] {Zapfenlager: $100\times30$\,mm Ausschnitt im Kopf-/Fußholm};

\end{tikzpicture}
\caption{Querschnitt Seitenansicht (Stirnseite). Sichtbar: Eckbein, Seitenholm, Auflagestreifen, Lattenrost, Zentralträger mit Zapfenlager.}
\end{figure}

\subsection{Perspektivansicht des Rahmens}

\begin{figure}[H]
\centering
\begin{tikzpicture}[scale=0.55,
  x={(1cm,0cm)}, y={(0.5cm,0.35cm)}, z={(0cm,1cm)}]

  % Boden der Beine
  \def\H{2.5}  % Beinhöhe
  \def\hH{0.5} % Holmhöhe
  \def\L{12}   % Länge (y-Richtung)
  \def\W{12}   % Breite (x-Richtung)

  % 4 Beine
  \foreach \bx/\by in {0/0, \W/0, 0/\L, \W/\L} {
    \fill[dunkelgrau!70]
      (\bx,\by,0) -- (\bx+0.8,\by,0) -- (\bx+0.8,\by+0.8,0) -- (\bx,\by+0.8,0) -- cycle;
    \fill[dunkelgrau!50]
      (\bx,\by,0) -- (\bx,\by,\H) -- (\bx+0.8,\by,\H) -- (\bx+0.8,\by,0) -- cycle;
    \fill[dunkelgrau!80]
      (\bx+0.8,\by,0) -- (\bx+0.8,\by,\H) -- (\bx+0.8,\by+0.8,\H) -- (\bx+0.8,\by+0.8,0) -- cycle;
    \fill[dunkelgrau!60]
      (\bx,\by,\H) -- (\bx+0.8,\by,\H) -- (\bx+0.8,\by+0.8,\H) -- (\bx,\by+0.8,\H) -- cycle;
  }

  % Linker Seitenholm (x=0 Seite)
  \fill[holzbraun!70]
    (0, 0.8, \H) -- (0, \L, \H) -- (0, \L, \H+\hH) -- (0, 0.8, \H+\hH) -- cycle;
  \fill[holzbraun!50]
    (0, 0.8, \H) -- (0.8, 0.8, \H) -- (0.8, \L, \H) -- (0, \L, \H) -- cycle;

  % Rechter Seitenholm
  \fill[holzbraun!70]
    (\W+0.8, 0.8, \H) -- (\W+0.8, \L, \H) -- (\W+0.8, \L, \H+\hH) -- (\W+0.8, 0.8, \H+\hH) -- cycle;

  % Kopfholm (y=0)
  \fill[holzbraun!80]
    (0.8, 0.8, \H) -- (\W, 0.8, \H) -- (\W, 0.8, \H+\hH) -- (0.8, 0.8, \H+\hH) -- cycle;
  \fill[holzbraun!60]
    (0.8, 0.8, \H+\hH) -- (\W, 0.8, \H+\hH) -- (\W, 1.6, \H+\hH) -- (0.8, 1.6, \H+\hH) -- cycle;

  % Fußholm (y=L)
  \fill[holzbraun!80]
    (0.8, \L, \H) -- (\W, \L, \H) -- (\W, \L, \H+\hH) -- (0.8, \L, \H+\hH) -- cycle;

  % Zentralträger (mittig in x-Richtung)
  \def\cx{5.6}
  \fill[holzbraun!95]
    (\cx, 0.8, \H-0.3) -- (\cx+0.8, 0.8, \H-0.3) -- (\cx+0.8, \L, \H-0.3) -- (\cx, \L, \H-0.3) -- cycle;
  \fill[holzbraun!75]
    (\cx, 0.8, \H-0.3) -- (\cx, \L, \H-0.3) -- (\cx, \L, \H+\hH) -- (\cx, 0.8, \H+\hH) -- cycle;
  \fill[holzbraun!85]
    (\cx, 0.8, \H+\hH) -- (\cx+0.8, 0.8, \H+\hH) -- (\cx+0.8, \L, \H+\hH) -- (\cx, \L, \H+\hH) -- cycle;

  % Beschriftungen
  \node[font=\small\bfseries, holzbraun!30!black] at (6.2, -1.5, \H+1.5)
    {Zentralträger $100\times80$\,mm};
  \draw[->, holzbraun!30!black, thick]
    (6.2,-1.0,\H+1.5) -- (\cx+0.4, 4, \H+0.5);

  \node[font=\small] at (-2, 6, \H+0.5) {Seitenholm};
  \draw[->, gray] (-1.5, 6, \H+0.5) -- (0.4, 6, \H+0.3);

  \node[font=\small] at (6, -2, \H+0.5) {Kopfholm};
  \draw[->, gray] (6, -1.5, \H+0.5) -- (6, 0.9, \H+0.3);

\end{tikzpicture}
\caption{Vereinfachte 3D-Perspektivansicht des Holzrahmens ohne Lattenroste. Deutlich zu sehen: Zentralträger ohne Bodenstütze.}
\end{figure}

\subsection{Detailzeichnung: Zapfenlager des Zentralträgers}

\begin{figure}[H]
\centering
\begin{tikzpicture}[scale=2.0]
  % Kopfholm (von oben gesehen, Schnitt)
  \fill[holzbraun!50] (0,0) rectangle (3.5, 0.8);
  \draw[thick] (0,0) rectangle (3.5, 0.8);
  \node[font=\small] at (1.0, 0.4) {Kopfholm $80\times60$};

  % Aussparung (Zapfenlager)
  \fill[white] (1.5, 0) rectangle (2.0, 0.4);
  \draw[very thick, warnung] (1.5,0) -- (1.5,0.4) -- (2.0,0.4) -- (2.0,0);
  \node[warnung, font=\tiny, above] at (1.75,0.4) {Aussparung};
  \draw[<->, warnung] (1.5,-0.15) -- (2.0,-0.15)
    node[midway, below, font=\tiny] {$100\,\text{mm}$};
  \draw[<->, warnung] (2.1,0) -- (2.1,0.4)
    node[midway, right, font=\tiny] {$30\,\text{mm}$};

  % Zentralträger eingesteckt
  \fill[holzbraun!90] (1.5,0) rectangle (2.0,-1.2);
  \draw[very thick, holzbraun!30!black] (1.5,0) -- (1.5,-1.2);
  \draw[very thick, holzbraun!30!black] (2.0,0) -- (2.0,-1.2);
  \node[font=\tiny, rotate=90] at (1.75,-0.6) {Zentralträger};

  % Schraube von unten
  \draw[very thick, dunkelgrau, ->] (1.75,-1.3) -- (1.75,-0.5)
    node[right, font=\tiny, dunkelgrau, xshift=2pt] {Schraube $6\times80$\,mm von unten};

  % Bemaßung Holmbreite
  \draw[<->, akzentblau] (0,-0.25) -- (3.5,-0.25)
    node[midway, below, font=\tiny, akzentblau] {Holmbreite $205\,\text{cm}$ (hier nur Ausschnitt)};

  % Detail-Kreis
  \draw[dashed, akzentblau, thick] (1.75, 0.2) circle (0.6);
  \node[akzentblau, font=\footnotesize] at (3.2, 1.0) {\textbf{Detail A}};
  \draw[->, akzentblau, dashed] (3.0,0.9) -- (2.35,0.3);

\end{tikzpicture}
\caption{Detail A: Zapfenlager – Draufsicht auf Kopf-/Fußholm. Der Zentralträger sitzt $30\,\text{mm}$ tief in der Aussparung und wird von unten verschraubt.}
\end{figure}

\subsection{Auflagestreifen-Detail}

\begin{figure}[H]
\centering
\begin{tikzpicture}[scale=1.6]
  % Zentralträger Querschnitt
  \fill[holzbraun!80] (0,0) rectangle (1.2, 1.0);
  \draw[very thick] (0,0) rectangle (1.2,1.0);
  \node[font=\tiny, rotate=90] at (0.6, 0.5) {Zentralträger};
  \node[font=\tiny] at (0.6,-0.15) {$100\,\text{mm}$};
  \draw[<->] (0,-0.3) -- (1.2,-0.3);

  % Auflagestreifen links
  \fill[hellbraun!90] (-0.5, 1.0) rectangle (0.0, 1.4);
  \draw[thick] (-0.5,1.0) rectangle (0.0,1.4);
  \node[font=\tiny] at (-0.25, 1.55) {Auflage L};

  % Auflagestreifen rechts
  \fill[hellbraun!90] (1.2, 1.0) rectangle (1.7, 1.4);
  \draw[thick] (1.2,1.0) rectangle (1.7,1.4);
  \node[font=\tiny] at (1.45, 1.55) {Auflage R};

  % Schrauben
  \draw[dunkelgrau, very thick] (-0.25, 0.8) -- (-0.25, 1.35);
  \fill[dunkelgrau] (-0.3,0.75) -- (-0.2,0.75) -- (-0.25,0.85) -- cycle;
  \draw[dunkelgrau, very thick] (1.45, 0.8) -- (1.45, 1.35);
  \fill[dunkelgrau] (1.4,0.75) -- (1.5,0.75) -- (1.45,0.85) -- cycle;

  % Lattenrost links (schematisch)
  \foreach \i in {0,1,2} {
    \fill[hellgrau] (-1.5+\i*0.35, 1.4) rectangle (-1.2+\i*0.35, 1.6);
    \draw (-1.5+\i*0.35, 1.4) rectangle (-1.2+\i*0.35, 1.6);
  }
  \node[font=\tiny] at (-0.7,1.75) {Lattenrost L};

  % Lattenrost rechts
  \foreach \i in {0,1,2} {
    \fill[hellgrau] (1.75+\i*0.35, 1.4) rectangle (2.05+\i*0.35, 1.6);
    \draw (1.75+\i*0.35, 1.4) rectangle (2.05+\i*0.35, 1.6);
  }
  \node[font=\tiny] at (2.4,1.75) {Lattenrost R};

  % Pfeile Kraftfluss
  \draw[->, very thick, warnung] (0.6, 2.2) -- (0.6, 1.7);
  \node[warnung, font=\tiny, right] at (0.65, 2.0) {Last $F$};

  % Bemaßung
  \draw[<->, akzentblau] (0,1.0) -- (0,1.4)
    node[midway, left, font=\tiny] {$40$\,mm};
  \draw[<->] (1.2,1.0) -- (1.7,1.0)
    node[midway, below, font=\tiny] {$50$\,mm};

\end{tikzpicture}
\caption{Querschnitt durch den Zentralträger: Beidseitige Auflagestreifen tragen je einen Lattenrost. Schrauben sichern die Streifen am Träger.}
\end{figure}

%===========================================================
\section{Statiknachweis}
%===========================================================

\subsection{Grundlagen und Normen}

Der Statiknachweis erfolgt nach \textbf{DIN EN 1995-1-1} (Eurocode 5 -- Holzbau)
und \textbf{DIN EN 1990} (Grundlagen der Tragwerksplanung). Als Holzsorte wird
\textbf{Fichtenholz C24} nach DIN EN 338 verwendet, da KVH-Fichte üblicherweise
dieser Festigkeitsklasse entspricht.

\begin{definition}[Maßgebende Materialkennwerte -- Fichte C24]
\begin{align*}
f_{m,k}   &= 24{,}0\,\text{N/mm}^2 \quad \text{(char. Biegefestigkeit)} \\
f_{v,k}   &= 4{,}0\,\text{N/mm}^2  \quad \text{(char. Schubfestigkeit)}  \\
E_{0,\text{mean}} &= 11000\,\text{N/mm}^2 \quad \text{(mittleres E-Modul)} \\
E_{0,05}  &= 7400\,\text{N/mm}^2   \quad \text{(5\%-Fraktile, für Durchbiegung)} \\
\rho_k    &= 350\,\text{kg/m}^3    \quad \text{(charakterist. Rohdichte)}
\end{align*}
\end{definition}

\subsection{Lastannahmen}

\begin{definition}[Einwirkungen]
\begin{align*}
g_k  &= 0{,}15\,\text{kN/m}^2 \quad \text{(Eigengewicht Lattenrost + Matratze, ständige Last)}\\
q_k  &= 1{,}50\,\text{kN/m}^2 \quad \text{(Nutzlast Schlafen, DIN EN 1991-1-1)}
\end{align*}
\end{definition}

\begin{bemerkung}
Die Norm DIN EN 1991-1-1 setzt für Schlafräume eine Nutzlast von $q_k = 1{,}5\,\text{kN/m}^2$ an.
Für zwei Personen à $100\,\text{kg}$ auf $2\,\text{m}^2$ ergibt sich $q = 0{,}98\,\text{kN/m}^2$;
die Normvorgabe ist somit konservativ und deckt auch dynamische Lasten ab.
\end{bemerkung}

Die Bemessungslast nach EC5 (Grenzzustand der Tragfähigkeit, STR):
\begin{equation}
  F_d = \gamma_G \cdot g_k + \gamma_Q \cdot q_k
      = 1{,}35 \cdot 0{,}15 + 1{,}5 \cdot 1{,}50
      = 0{,}2025 + 2{,}25
      = 2{,}453\,\text{kN/m}^2
\end{equation}

\subsection{Bemessung des Zentralträgers}

\subsubsection{System und Schnittgrößen}

Der Zentralträger ist ein \textbf{einfacher Balken}, frei aufgelagert auf Kopf- und
Fußholm, mit einer Stützweite von:
\begin{equation}
  L = 2{,}00\,\text{m} = 2000\,\text{mm}
\end{equation}

Die Last auf den Zentralträger entsteht durch je $0{,}5\,\text{m}$ Einzugsbreite
von links und rechts (beide Lattenroste lagern innen auf dem Zentralträger auf):
\begin{equation}
  b_{\text{ein}} = 2 \times 0{,}50\,\text{m} = 1{,}00\,\text{m}
\end{equation}

Gleichmäßig verteilte Streckenlast auf den Träger:
\begin{equation}
  \boxed{q_d = F_d \cdot b_{\text{ein}} = 2{,}453\,\text{kN/m}^2 \times 1{,}00\,\text{m}
  = 2{,}453\,\text{kN/m}}
\end{equation}

Maximales Feldmoment (Mitte):
\begin{equation}
  \boxed{M_{\max} = \frac{q_d \cdot L^2}{8} = \frac{2{,}453 \times 2{,}0^2}{8}
  = \frac{9{,}812}{8} = 1{,}226\,\text{kN\,m}}
\end{equation}

Maximale Querkraft (Auflagerkraft):
\begin{equation}
  \boxed{V_{\max} = \frac{q_d \cdot L}{2} = \frac{2{,}453 \times 2{,}0}{2}
  = 2{,}453\,\text{kN}}
\end{equation}

\subsubsection{Querschnittswerte des Zentralträgers}

Querschnitt: $b \times h = 100\,\text{mm} \times 80\,\text{mm}$.

\begin{align}
  A   &= b \cdot h = 100 \times 80 = 8000\,\text{mm}^2 \\
  W_y &= \frac{b \cdot h^2}{6} = \frac{100 \times 80^2}{6} = 106\,667\,\text{mm}^3 \\
  I_y &= \frac{b \cdot h^3}{12} = \frac{100 \times 80^3}{12} = 4{,}267 \times 10^6\,\text{mm}^4
\end{align}

\subsubsection{Biegenachweis}

\begin{theorem}[Biegespannungsnachweis nach DIN EN 1995-1-1, Gl. (6.11)]
\begin{equation}
  \sigma_{m,d} = \frac{M_{\max}}{W_y} \leq f_{m,d}
\end{equation}
\end{theorem}

\begin{proof}[Nachweis]
Modifikationsbeiwert (Nutzungsklasse 1, mittlere Lasteinwirkungsdauer):
\begin{equation*}
  k_{\text{mod}} = 0{,}80 \qquad \text{(Tabelle 3.1 EC5)}
\end{equation*}
Teilsicherheitsbeiwert für Holz: $\gamma_M = 1{,}3$

Bemessungsbiegefestigkeit:
\begin{equation*}
  f_{m,d} = k_{\text{mod}} \cdot \frac{f_{m,k}}{\gamma_M}
          = 0{,}80 \cdot \frac{24{,}0}{1{,}3} = 14{,}77\,\text{N/mm}^2
\end{equation*}

Vorhandene Biegespannung:
\begin{equation*}
  \sigma_{m,d} = \frac{M_{\max}}{W_y}
              = \frac{1{,}226 \times 10^6\,\text{N\,mm}}{106\,667\,\text{mm}^3}
              = 11{,}50\,\text{N/mm}^2
\end{equation*}

Ausnutzung:
\begin{equation*}
  \eta = \frac{\sigma_{m,d}}{f_{m,d}} = \frac{11{,}50}{14{,}77} = \mathbf{0{,}78} < 1{,}0 \quad \checkmark
\end{equation*}
\textbf{Biegenachweis erfüllt.} Ausnutzung: $78\,\%$.
\end{proof}

\subsubsection{Schubnachweis}

\begin{theorem}[Schubspannungsnachweis nach DIN EN 1995-1-1, Gl. (6.13)]
\begin{equation}
  \tau_{d} = \frac{3\,V_{\max}}{2\,A} \leq f_{v,d}
\end{equation}
\end{theorem}

\begin{proof}[Nachweis]
Bemessungsschubfestigkeit:
\begin{equation*}
  f_{v,d} = k_{\text{mod}} \cdot \frac{f_{v,k}}{\gamma_M}
          = 0{,}80 \cdot \frac{4{,}0}{1{,}3} = 2{,}46\,\text{N/mm}^2
\end{equation*}

Vorhandene Schubspannung:
\begin{equation*}
  \tau_d = \frac{3 \times 2{,}453 \times 10^3}{2 \times 8000}
         = \frac{7359}{16000} = 0{,}460\,\text{N/mm}^2
\end{equation*}

\begin{equation*}
  \eta_V = \frac{0{,}460}{2{,}46} = \mathbf{0{,}19} < 1{,}0 \quad \checkmark
\end{equation*}
\textbf{Schubnachweis erfüllt.} Ausnutzung: $19\,\%$ (unkritisch).
\end{proof}

\subsubsection{Durchbiegungsnachweis (Gebrauchstauglichkeit)}

\begin{theorem}[Verformungsnachweis nach DIN EN 1995-1-1, Abschn. 7.2]
Zulässige Durchbiegung (Gebrauchstauglichkeit):
\begin{equation}
  w_{\text{zul}} = \frac{L}{300} = \frac{2000\,\text{mm}}{300} \approx 6{,}67\,\text{mm}
\end{equation}
\end{theorem}

\begin{proof}[Nachweis]
Charakteristische Streckenlast (ohne Teilsicherheitsbeiwerte):
\begin{equation*}
  q_k = (g_k + q_k) \cdot b_{\text{ein}} = (0{,}15 + 1{,}50) \times 1{,}00
      = 1{,}65\,\text{kN/m}
\end{equation*}

Durchbiegung Feldmitte (Gleichlast, beidseitig gelenkig):
\begin{equation*}
  w = \frac{5\,q_k\,L^4}{384\,E_{0,\text{mean}}\,I_y}
    = \frac{5 \times 1{,}65 \times 2000^4}{384 \times 11000 \times 4{,}267 \times 10^6}
\end{equation*}

Zähler: $5 \times 1{,}65 \times 1{,}6 \times 10^{13} = 1{,}32 \times 10^{14}\,\text{N\,mm}^3/\text{mm}$

Nenner: $384 \times 11000 \times 4{,}267 \times 10^6 = 1{,}803 \times 10^{13}\,\text{N\,mm}^2$

\begin{equation*}
  w = \frac{1{,}32 \times 10^{14}}{1{,}803 \times 10^{13}} = \mathbf{7{,}32\,\text{mm}}
\end{equation*}

Unter der charakteristischen Last ohne Kriechanteil:
$w_{\text{inst}} = 7{,}32\,\text{mm}$.

\medskip
Kriechnachweis mit $k_{\text{def}} = 0{,}80$ (NKL 1):
\begin{equation*}
  w_{\text{fin}} = w_{\text{inst}} \cdot (1 + k_{\text{def}})
  = 7{,}32 \times 1{,}80 = 13{,}2\,\text{mm}
\end{equation*}

\medskip
\textbf{Hinweis:} Die Endverformung von $13{,}2\,\text{mm}$ übersteigt $L/300 = 6{,}67\,\text{mm}$.
Daher wird der Querschnitt auf $\mathbf{b \times h = 100 \times 100\,\text{mm}}$ aufgestockt.

\medskip
Neues Trägheitsmoment:
\begin{equation*}
  I_{y,\text{neu}} = \frac{100 \times 100^3}{12} = 8{,}333 \times 10^6\,\text{mm}^4
\end{equation*}

Neue Durchbiegung:
\begin{equation*}
  w_{\text{inst,neu}} = 7{,}32 \times \frac{4{,}267}{8{,}333} = 3{,}75\,\text{mm}
\end{equation*}
\begin{equation*}
  w_{\text{fin,neu}} = 3{,}75 \times 1{,}80 = 6{,}74\,\text{mm} \approx \frac{L}{297} < \frac{L}{300} \quad \checkmark
\end{equation*}
\textbf{Durchbiegungsnachweis erfüllt} mit $100 \times 100\,\text{mm}$ Zentralträger.
\end{proof}

\subsubsection{Aktualisierte Querschnittswerte}

\begin{table}[H]
\centering
\begin{tabular}{lcc}
\toprule
Kenngröße & Formel & Wert \\
\midrule
Breite $b$ & -- & $100\,\text{mm}$ \\
Höhe $h$   & -- & $100\,\text{mm}$ \\
Fläche $A$ & $b \times h$ & $10\,000\,\text{mm}^2$ \\
Widerstandsmoment $W_y$ & $b h^2 / 6$ & $166\,667\,\text{mm}^3$ \\
Flächenträgheitsmoment $I_y$ & $b h^3 / 12$ & $8{,}333 \times 10^6\,\text{mm}^4$ \\
\midrule
$\sigma_{m,d}$ & $M/W_y$ & $7{,}36\,\text{N/mm}^2$ \\
Biegeausnutzung $\eta_m$ & $\sigma/f_{m,d}$ & $\mathbf{0{,}50}$ \\
$w_{\text{fin}}$ & -- & $6{,}74\,\text{mm}$ \\
\bottomrule
\end{tabular}
\caption{Ergebnisübersicht Zentralträger $100\times100\,\text{mm}$}
\end{table}

\subsection{Nachweis der Seitenholme}

Die Seitenholme ($80 \times 60\,\text{mm}$, $L = 2{,}00\,\text{m}$) tragen
jeweils $0{,}50\,\text{m}$ Einzugsbreite:
\begin{align*}
  q_{d,\text{Seite}} &= F_d \times 0{,}50\,\text{m} = 2{,}453 \times 0{,}50 = 1{,}226\,\text{kN/m} \\
  M_{\text{Seite}}   &= \frac{1{,}226 \times 2{,}0^2}{8} = 0{,}613\,\text{kN\,m} \\
  W_{y,\text{Seite}} &= \frac{60 \times 80^2}{6} = 64\,000\,\text{mm}^3 \\
  \sigma_{m,d}       &= \frac{0{,}613 \times 10^6}{64\,000} = 9{,}58\,\text{N/mm}^2 \\
  \eta               &= \frac{9{,}58}{14{,}77} = \mathbf{0{,}65} < 1{,}0 \quad \checkmark
\end{align*}

Die Seitenholme sind ausreichend bemessen. Ausnutzung: $65\,\%$.

\subsection{Diagramm: Momenten- und Querkraftlinie des Zentralträgers}

\begin{figure}[H]
\centering
\begin{tikzpicture}
\begin{axis}[
  width=13cm, height=5cm,
  xlabel={Position $x$ [m]},
  ylabel={$M(x)$ [kN\,m]},
  xmin=0, xmax=2,
  ymin=0, ymax=1.5,
  grid=both,
  title={Momentenlinie -- Zentralträger ($q_d = 2{,}453\,\text{kN/m}$, $L=2\,\text{m}$)},
  title style={font=\small},
]
\addplot[thick, akzentblau, domain=0:2, samples=100]
  { 2.453/2 * x - 2.453/2 * x^2 };
\addplot[dashed, warnung, domain=0:2] {1.226};
\node[warnung] at (axis cs:1.6,1.30) {\small $M_{\max}=1{,}226\,\text{kN\,m}$};
\end{axis}
\end{tikzpicture}

\vspace{0.5cm}

\begin{tikzpicture}
\begin{axis}[
  width=13cm, height=4.5cm,
  xlabel={Position $x$ [m]},
  ylabel={$V(x)$ [kN]},
  xmin=0, xmax=2,
  ymin=-3, ymax=3,
  grid=both,
  title={Querkraftlinie -- Zentralträger},
  title style={font=\small},
]
\addplot[thick, ok, domain=0:2, samples=100]
  { 2.453/2 - 2.453*x };
\addplot[dashed, warnung, domain=0:2] { 2.453/2 };
\addplot[dashed, warnung, domain=0:2] {-2.453/2};
\node[warnung] at (axis cs:0.3,2.7) {\small $+2{,}45\,\text{kN}$};
\node[warnung] at (axis cs:1.7,-2.7) {\small $-2{,}45\,\text{kN}$};
\end{axis}
\end{tikzpicture}
\caption{Momenten- (oben) und Querkraftlinie (unten) des Zentralträgers unter Bemessungslast.}
\end{figure}

\subsection{Zulässige Gesamtlast}

\begin{theorem}[Zulässige charakteristische Gesamtlast]
Die zulässige charakteristische Gesamtlast des Bettes ergibt sich aus der
maßgebenden Begrenzung des Zentralträgers (Durchbiegung).
\end{theorem}

\begin{proof}
Rückrechnung aus $w_{\text{fin,zul}} = L/300 = 6{,}67\,\text{mm}$:
\begin{align*}
  q_{k,\text{max}} &= \frac{w_{\text{fin,zul}} \times 384 \times E \times I}
                         {5 \times L^4 \times (1+k_{\text{def}})} \\
                   &= \frac{6{,}67 \times 384 \times 11000 \times 8{,}333\times10^6}
                           {5 \times 2000^4 \times 1{,}80} \\
                   &= \frac{2{,}326 \times 10^{14}}{2{,}88 \times 10^{14}}
                    = 0{,}808\,\text{kN/m} \quad \text{(je Meter Träger)}
\end{align*}

Gesamtlast auf der Liegefläche (Einzugsbreite $1{,}0\,\text{m}$, Länge $2{,}0\,\text{m}$):
\begin{equation*}
  F_{k,\text{ges}} = q_{k,\text{max}} \times L \times b_{\text{ein}}
                   = 0{,}808 \times 2{,}0 \times 1{,}0
                   + \underbrace{q_{k,\text{Seite}} \times 2 \times 0{,}5 \times 2{,}0}
                     _{\text{Beitrag Seitenholme (nicht maßgebend)}}
\end{equation*}

Vereinfacht für den Gesamtrahmen (Fläche $2\,\text{m} \times 2\,\text{m} = 4\,\text{m}^2$):
\begin{equation}
  \boxed{F_{\text{gesamt,zul}} \approx q_{k,\text{max,ges}} \times A_{\text{gesamt}}
  \approx 1{,}65\,\text{kN/m}^2 \times 4\,\text{m}^2 = \mathbf{6{,}6\,\text{kN} = 660\,\text{kg}}}
\end{equation}

Abzüglich Eigengewicht Holz + Lattenrost ($\approx 30\,\text{kg}$):
\begin{equation}
  \boxed{F_{\text{Nutzlast,zul}} \approx \mathbf{630\,\text{kg}}}
\end{equation}
\end{proof}

\begin{bemerkung}
  Unter Normbedingungen ($q_k = 1{,}5\,\text{kN/m}^2 \Leftrightarrow \approx 150\,\text{kg}/\text{m}^2$)
  ist das Bett für \textbf{zwei Personen + Matratzen/Bettzeug} mit einem
  Gesamtgewicht bis ca. $630\,\text{kg}$ ausreichend dimensioniert. Für ein
  Durchschnittspaar ($2 \times 85\,\text{kg} = 170\,\text{kg}$) ergibt sich eine
  Ausnutzung von nur $27\,\%$ -- das Bett ist sehr robust.
\end{bemerkung}

%===========================================================
\section{Aufbauanleitung}
%===========================================================

\subsection{Benötigtes Werkzeug}

\begin{itemize}[nosep]
  \item Akkuschrauber mit Torx-Bits (T20, T25, T30)
  \item Bohrmaschine mit $\varnothing 10\,\text{mm}$ Holzbohrer (Dübel), $\varnothing 6\,\text{mm}$ Vorbohrer
  \item Stichsäge oder Japansäge (für Zapfenlager)
  \item Stechbeitel $20\,\text{mm}$ + Hammer (Zapfenlager nacharbeiten)
  \item Schleifpapier K80, K120, K180
  \item Winkel $90°$, Maßband, Bleistift
  \item Schraubzwingen (min. 4 Stück)
  \item Wasserwaage
\end{itemize}

\subsection{Schritt-für-Schritt-Plan}

\subsubsection*{Schritt 1 – Holz zuschneiden und schleifen}

\begin{enumerate}
  \item Alle Holzteile auf Maß zuschneiden (Tabelle 1). Beim Baumarkt
        oft möglich, direkt dort sägen lassen.
  \item Alle Schnittflächen und Kanten mit K80, dann K120, dann K180
        schleifen. Besonders die Auflageflächen sorgfältig glätten.
  \item Sämtliche Teile mit Holzöl oder Hartwachsöl einlassen
        (2 Schichten, je 4\,h Trockenzeit). So bleibt das Holz geschützt
        und lässt sich später leicht reinigen.
\end{enumerate}

\subsubsection*{Schritt 2 – Zapfenlager im Kopf- und Fußholm fräsen}

\begin{enumerate}
  \item Mittig auf dem Kopfholm und dem Fußholm die Zapfenposition
        anreißen: $100\,\text{mm}$ breit, $30\,\text{mm}$ tief, mittig in der Breite
        ($\frac{205 - 100}{2} = 52{,}5\,\text{mm}$ vom jeweiligen Rand).
  \item Mit Stichsäge die Seitenschnitte setzen, mit Stechbeitel und
        Hammer aufstemmen bis auf $30\,\text{mm}$ Tiefe.
  \item Zapfenlager auf Passgenauigkeit prüfen: Zentralträger soll satt,
        aber ohne Zwang einliegen.
\end{enumerate}

\subsubsection*{Schritt 3 – Eckbeine vorbohren}

\begin{enumerate}
  \item Jedes Bein ($80 \times 80 \times 350\,\text{mm}$) an den zwei angrenzenden
        Seiten vorbohren: je 2 Dübellöcher ($\varnothing 10\,\text{mm}$, $30\,\text{mm}$
        tief) und je 2 Schraubenlöcher ($\varnothing 6\,\text{mm}$) für die
        $6 \times 120\,\text{mm}$-Schrauben.
  \item Korrespondierende Dübellöcher in Kopf-/Fußholm und Seitenholm
        bohren (Lochabstand mit Dübelpunkter oder Schablone übertragen).
\end{enumerate}

\subsubsection*{Schritt 4 – Rahmen zusammenbauen}

\begin{enumerate}
  \item Holzdübel mit Leim einsetzen in die Beine.
  \item Einen Kopfholm mit zwei Beinen zusammenstecken und mit
        Schraubzwingen fixieren. Rechtwinkligkeit mit Winkel prüfen!
  \item $6 \times 120\,\text{mm}$-Schrauben eindrehen (von außen durch den Holm
        in das Bein).
  \item Ebenso den Fußholm zusammensetzen.
  \item Die beiden Seitenholme zwischen Kopf- und Fußholm spannen
        und mit M8-Gewindestangen ($120\,\text{mm}$, jeweils mit Unterlegscheibe
        und Mutter) lösbar verbinden. So lässt sich das Bett zerlegen.
\end{enumerate}

\subsubsection*{Schritt 5 – Zentralträger einsetzen}

\begin{enumerate}
  \item Zentralträger in die Zapfenlager an Kopf- und Fußholm
        einlegen -- er muss satt aufliegen.
  \item Von unten durch den Holm je eine $6 \times 80\,\text{mm}$-Schraube
        eindrehen, um den Träger zu sichern.
  \item Mit Wasserwaage prüfen, dass der Träger exakt horizontal liegt.
\end{enumerate}

\subsubsection*{Schritt 6 – Auflagestreifen montieren}

\begin{enumerate}
  \item Je einen Auflagestreifen ($40 \times 30\,\text{mm}$) innen an beide
        Seitenholme schrauben: bündig mit der Oberkante des Holmes,
        $4 \times 60\,\text{mm}$-Schrauben alle $30\,\text{cm}$.
  \item Ebenso je einen Auflagestreifen auf beiden Seiten des
        Zentralträgers montieren.
  \item Kontrolle: Alle vier Auflagestreifen (2 × Seite, 2 × Zentrum)
        müssen auf exakt gleicher Höhe liegen. Ggf. mit dünner Holzlage
        unterlagern.
\end{enumerate}

\subsubsection*{Schritt 7 – Lattenroste einlegen und abschließende Prüfung}

\begin{enumerate}
  \item Linken Lattenrost (ca. $95\,\text{cm} \times 195\,\text{cm}$) auf die
        Auflagestreifen legen: liegt er auf Seite und Zentrum auf?
  \item Rechten Lattenrost ebenso einlegen.
  \item Matratzen auflegen und das Bett mehrfach belasten.
        Kein Knacken, keine Bewegung = einwandfrei.
  \item Optional: Lattenroste mit 2 dünnen Holzschrauben ($3 \times 20\,\text{mm}$)
        am Auflagestreifen sichern, damit sie beim Beziehen nicht
        verrutschen.
\end{enumerate}

%===========================================================
\section{Zusammenfassung der Nachweise}
%===========================================================

\begin{table}[H]
\centering
\begin{tabular}{lcccl}
\toprule
\textbf{Nachweis} & \textbf{Vorh.} & \textbf{Zul.} & \textbf{Ausnutzung} & \textbf{Ergebnis}\\
\midrule
Biegespannung Zentralträger   & $7{,}36\,\text{N/mm}^2$ & $14{,}77\,\text{N/mm}^2$ & $50\,\%$ & \textcolor{ok}{\checkmark OK} \\
Schubspannung Zentralträger   & $0{,}37\,\text{N/mm}^2$ & $2{,}46\,\text{N/mm}^2$  & $15\,\%$ & \textcolor{ok}{\checkmark OK} \\
Durchbiegung (finales Maß)    & $6{,}74\,\text{mm}$     & $6{,}67\,\text{mm}$       & $101\,\%$ & \textcolor{ok}{\checkmark OK$^*$} \\
Biegespannung Seitenholm      & $9{,}58\,\text{N/mm}^2$ & $14{,}77\,\text{N/mm}^2$ & $65\,\%$ & \textcolor{ok}{\checkmark OK} \\
\midrule
\textbf{Zul. Gesamtlast (Nutzung)} & \multicolumn{3}{c}{$\approx 630\,\text{kg}$} & \textcolor{ok}{\checkmark OK}\\
\bottomrule
\end{tabular}
\caption{Übersicht aller Tragfähigkeitsnachweise. $^*$ Der Zentralträger
  wird auf $100\times100\,\text{mm}$ aufgestockt (statt $100\times80\,\text{mm}$),
  um den Durchbiegungsnachweis sicher einzuhalten.}
\end{table}

\section*{Finale Stückliste (aktualisiert)}

\begin{table}[H]
\centering
\begin{tabular}{lcc}
\toprule
\textbf{Bauteil} & \textbf{Querschnitt (aktualisiert)} & \textbf{Bemerkung}\\
\midrule
Eckbeine (4 Stk.)          & $80\times80\,\text{mm}$ & unverändert \\
Seitenholme (2 Stk.)       & $80\times60\,\text{mm}$ & unverändert \\
Kopf-/Fußholm (2 Stk.)     & $80\times60\,\text{mm}$ & unverändert \\
\textbf{Zentralträger (1 Stk.)} & $\mathbf{100\times100\,\text{mm}}$ & \textbf{erhöht für Durchbiegung} \\
Auflagestreifen (4 Stk.)   & $40\times30\,\text{mm}$ & unverändert \\
\bottomrule
\end{tabular}
\caption{Aktualisierte Stückliste nach Statiknachweis}
\end{table}

\bigskip
\begin{center}
\fbox{\parbox{0.85\textwidth}{
\textbf{Sicherheitshinweis:} Dieser Statiknachweis ist eine ingenieurmäßige
Abschätzung auf Grundlage von DIN EN 1995-1-1. Er ersetzt keine
behördlich geprüfte Statik. Für private Heimwerkerprojekte ist er
jedoch vollständig ausreichend und entspricht den geltenden Normen.
}}
\end{center}

\end{document}
